% Figura 2 -- Ricart-Agrawala com tres participantes.
% Uso: % Figura 2 -- Ricart-Agrawala com tres participantes.
% Uso: % Figura 2 -- Ricart-Agrawala com tres participantes.
% Uso: % Figura 2 -- Ricart-Agrawala com tres participantes.
% Uso: \input{figuras/sequencia-ra} dentro de um ambiente figure.
\begin{tikzpicture}[
    font=\footnotesize,
    cab/.style={draw, rounded corners=2pt, minimum width=1.15cm, minimum height=0.6cm, fill=black!6},
    seta/.style={-{Latex[length=1.7mm]}, semithick},
    rot/.style={font=\scriptsize, inner sep=1.5pt, fill=white}
  ]

  % Cabecalhos e linhas de vida
  \node[cab] (h1) at (0,0)   {N1};
  \node[cab] (h2) at (3.2,0) {N2};
  \node[cab] (h3) at (6.4,0) {N3};
  \foreach \x in {0,3.2,6.4} { \draw[dashed, gray] (\x,-0.32) -- (\x,-4.45); }

  % REQUEST do N1 para os dois peers
  \draw[seta] (0,-1.00) -- node[rot, above] {REQUEST $(ts=5)$} (3.2,-1.00);
  \draw[seta] (0,-1.55) -- node[rot, above] {REQUEST $(ts=5)$} (6.4,-1.55);

  % Permissoes
  \draw[seta] (3.2,-2.10) -- node[rot, above] {REPLY} (0,-2.10);
  \draw[seta] (6.4,-2.55) -- node[rot, above] {REPLY} (0,-2.55);

  % N1 entra na secao critica com as duas permissoes
  \fill[black!12] (-0.42,-3.00) rectangle (0.42,-4.15);
  \node[font=\scriptsize, align=center, text width=1.05cm] at (0,-3.58)
    {se\c{c}\~ao\\cr\'itica};

\end{tikzpicture}
 dentro de um ambiente figure.
\begin{tikzpicture}[
    font=\footnotesize,
    cab/.style={draw, rounded corners=2pt, minimum width=1.15cm, minimum height=0.6cm, fill=black!6},
    seta/.style={-{Latex[length=1.7mm]}, semithick},
    rot/.style={font=\scriptsize, inner sep=1.5pt, fill=white}
  ]

  % Cabecalhos e linhas de vida
  \node[cab] (h1) at (0,0)   {N1};
  \node[cab] (h2) at (3.2,0) {N2};
  \node[cab] (h3) at (6.4,0) {N3};
  \foreach \x in {0,3.2,6.4} { \draw[dashed, gray] (\x,-0.32) -- (\x,-4.45); }

  % REQUEST do N1 para os dois peers
  \draw[seta] (0,-1.00) -- node[rot, above] {REQUEST $(ts=5)$} (3.2,-1.00);
  \draw[seta] (0,-1.55) -- node[rot, above] {REQUEST $(ts=5)$} (6.4,-1.55);

  % Permissoes
  \draw[seta] (3.2,-2.10) -- node[rot, above] {REPLY} (0,-2.10);
  \draw[seta] (6.4,-2.55) -- node[rot, above] {REPLY} (0,-2.55);

  % N1 entra na secao critica com as duas permissoes
  \fill[black!12] (-0.42,-3.00) rectangle (0.42,-4.15);
  \node[font=\scriptsize, align=center, text width=1.05cm] at (0,-3.58)
    {se\c{c}\~ao\\cr\'itica};

\end{tikzpicture}
 dentro de um ambiente figure.
\begin{tikzpicture}[
    font=\footnotesize,
    cab/.style={draw, rounded corners=2pt, minimum width=1.15cm, minimum height=0.6cm, fill=black!6},
    seta/.style={-{Latex[length=1.7mm]}, semithick},
    rot/.style={font=\scriptsize, inner sep=1.5pt, fill=white}
  ]

  % Cabecalhos e linhas de vida
  \node[cab] (h1) at (0,0)   {N1};
  \node[cab] (h2) at (3.2,0) {N2};
  \node[cab] (h3) at (6.4,0) {N3};
  \foreach \x in {0,3.2,6.4} { \draw[dashed, gray] (\x,-0.32) -- (\x,-4.45); }

  % REQUEST do N1 para os dois peers
  \draw[seta] (0,-1.00) -- node[rot, above] {REQUEST $(ts=5)$} (3.2,-1.00);
  \draw[seta] (0,-1.55) -- node[rot, above] {REQUEST $(ts=5)$} (6.4,-1.55);

  % Permissoes
  \draw[seta] (3.2,-2.10) -- node[rot, above] {REPLY} (0,-2.10);
  \draw[seta] (6.4,-2.55) -- node[rot, above] {REPLY} (0,-2.55);

  % N1 entra na secao critica com as duas permissoes
  \fill[black!12] (-0.42,-3.00) rectangle (0.42,-4.15);
  \node[font=\scriptsize, align=center, text width=1.05cm] at (0,-3.58)
    {se\c{c}\~ao\\cr\'itica};

\end{tikzpicture}
 dentro de um ambiente figure.
\begin{tikzpicture}[
    font=\footnotesize,
    cab/.style={draw, rounded corners=2pt, minimum width=1.15cm, minimum height=0.6cm, fill=black!6},
    seta/.style={-{Latex[length=1.7mm]}, semithick},
    rot/.style={font=\scriptsize, inner sep=1.5pt, fill=white}
  ]

  % Cabecalhos e linhas de vida
  \node[cab] (h1) at (0,0)   {N1};
  \node[cab] (h2) at (3.2,0) {N2};
  \node[cab] (h3) at (6.4,0) {N3};
  \foreach \x in {0,3.2,6.4} { \draw[dashed, gray] (\x,-0.32) -- (\x,-4.45); }

  % REQUEST do N1 para os dois peers
  \draw[seta] (0,-1.00) -- node[rot, above] {REQUEST $(ts=5)$} (3.2,-1.00);
  \draw[seta] (0,-1.55) -- node[rot, above] {REQUEST $(ts=5)$} (6.4,-1.55);

  % Permissoes
  \draw[seta] (3.2,-2.10) -- node[rot, above] {REPLY} (0,-2.10);
  \draw[seta] (6.4,-2.55) -- node[rot, above] {REPLY} (0,-2.55);

  % N1 entra na secao critica com as duas permissoes
  \fill[black!12] (-0.42,-3.00) rectangle (0.42,-4.15);
  \node[font=\scriptsize, align=center, text width=1.05cm] at (0,-3.58)
    {se\c{c}\~ao\\cr\'itica};

\end{tikzpicture}
